\begin{codebox}
\codeline{\textcolor{cmtgray}{\#\ 全书知识地图与学习路径}}
\codeline{\textcolor{cmtgray}{\#\ Book\ Roadmap\ and\ Learning\ Paths}}
\codeline{\textcolor{cmtgray}{\#}}
\codeline{\textcolor{cmtgray}{\#\ 全书九章构成"理论\ \(\rightarrow\)\ 方法\ \(\rightarrow\)\ 语言\ \(\rightarrow\)\ 推理\ \(\rightarrow\)\ 应用\ \(\rightarrow\)\ 融合\ \(\rightarrow\)\ 实战"的完整链路}}
\codeblank{}
\codeline{\textcolor{cmtgray}{\#\ ==============================}}
\codeline{\textcolor{cmtgray}{\#\ 1.\ 全书结构}}
\codeline{\textcolor{cmtgray}{\#\ ==============================}}
\codeline{\textcolor{cmtgray}{\#}}
\codeline{\textcolor{cmtgray}{\#\ 第1章\ 绪论}}
\codeline{\textcolor{cmtgray}{\#\ \ \ └──\ 哲学本体论\ \(\rightarrow\)\ 工程本体论的转向；AI时代定位}}
\codeline{\textcolor{cmtgray}{\#}}
\codeline{\textcolor{cmtgray}{\#\ 第2章\ 本体论基础理论\ \ \ \ \ \ \ \ \ \ ←——\ 理论地基}}
\codeline{\textcolor{cmtgray}{\#\ \ \ └──\ 类/属性/公理/实例；描述逻辑；一阶逻辑；单调与非单调推理}}
\codeline{\textcolor{cmtgray}{\#}}
\codeline{\textcolor{cmtgray}{\#\ 第3章\ 本体构建方法论\ \ \ \ \ \ \ \ \ \ ←——\ 怎么设计}}
\codeline{\textcolor{cmtgray}{\#\ \ \ └──\ 能力问题(CQ)；Ontology\ 101；METHONTOLOGY；OntoClean评估}}
\codeline{\textcolor{cmtgray}{\#}}
\codeline{\textcolor{cmtgray}{\#\ 第4章\ 本体描述语言与工具\ \ \ \ \ \ ←——\ 怎么表达}}
\codeline{\textcolor{cmtgray}{\#\ \ \ └──\ RDF/RDFS/OWL；SPARQL；Protégé}}
\codeline{\textcolor{cmtgray}{\#}}
\codeline{\textcolor{cmtgray}{\#\ 第5章\ 知识表示与推理机制\ \ \ \ \ \ ←——\ 怎么推理}}
\codeline{\textcolor{cmtgray}{\#\ \ \ └──\ Tableau算法；SWRL规则；时态推理；概率推理}}
\codeline{\textcolor{cmtgray}{\#}}
\codeline{\textcolor{cmtgray}{\#\ 第6章\ 工程领域应用案例\ \ \ \ \ \ \ \ ←——\ 用在哪里}}
\codeline{\textcolor{cmtgray}{\#\ \ \ └──\ 制造调度；自动驾驶；BIM；航空FMEA}}
\codeline{\textcolor{cmtgray}{\#}}
\codeline{\textcolor{cmtgray}{\#\ 第7章\ 知识图谱构建与管理\ \ \ \ \ \ ←——\ 怎么规模化}}
\codeline{\textcolor{cmtgray}{\#\ \ \ └──\ 抽取/对齐/融合；存储选型；SHACL质量校验}}
\codeline{\textcolor{cmtgray}{\#}}
\codeline{\textcolor{cmtgray}{\#\ 第8章\ 本体与大语言模型\ \ \ \ \ \ \ \ ←——\ 怎么与AI融合}}
\codeline{\textcolor{cmtgray}{\#\ \ \ └──\ 幻觉控制；GraphRAG；Text2SPARQL；本体引导Agent}}
\codeline{\textcolor{cmtgray}{\#}}
\codeline{\textcolor{cmtgray}{\#\ 第9章\ 综合实战案例\ \ \ \ \ \ \ \ \ \ \ \ ←——\ 全栈落地}}
\codeline{\textcolor{cmtgray}{\#\ \ \ └──\ 智能制造知识管理系统（OWL\ +\ Jena\ +\ owlready2）}}
\codeblank{}
\codeline{\textcolor{cmtgray}{\#\ ==============================}}
\codeline{\textcolor{cmtgray}{\#\ 2.\ 章节依赖关系}}
\codeline{\textcolor{cmtgray}{\#\ ==============================}}
\codeline{\textcolor{cmtgray}{\#\ 第2章\ ──\(\rightarrow\)\ 第3章\ ──\(\rightarrow\)\ 第4章\ ──\(\rightarrow\)\ 第5章}}
\codeline{\textcolor{cmtgray}{\#\ \ \ │\ \ \ \ \ \ \ \ \ \ \ \ \ \ \ \ \ \ \ \ │\ \ \ \ \ \ \ \ \ \ │}}
\codeline{\textcolor{cmtgray}{\#\ \ \ └────────────────────┴──\(\rightarrow\)\ 第6章\ ┘}}
\codeline{\textcolor{cmtgray}{\#\ \ \ \ \ \ \ \ \ \ \ \ \ \ \ \ \ \ \ \ \ \ \ \ \ \ \ \ \ \ │}}
\codeline{\textcolor{cmtgray}{\#\ \ \ \ \ \ \ \ \ \ \ \ \ \ \ \ \ \ \ \ 第7章\ ←──┴──\(\rightarrow\)\ 第8章}}
\codeline{\textcolor{cmtgray}{\#\ \ \ \ \ \ \ \ \ \ \ \ \ \ \ \ \ \ \ \ \ \ └────┬────┘}}
\codeline{\textcolor{cmtgray}{\#\ \ \ \ \ \ \ \ \ \ \ \ \ \ \ \ \ \ \ \ \ \ \ \ \ \ 第9章}}
\codeblank{}
\codeline{\textcolor{cmtgray}{\#\ 最短可用路径（最快做出能跑的系统）：}}
\codeline{\textcolor{cmtgray}{\#\ \ \ 第4章（语言）\(\rightarrow\)\ 第7章（入库）\(\rightarrow\)\ 第9章（实战）}}
\codeblank{}
\codeline{\textcolor{cmtgray}{\#\ ==============================}}
\codeline{\textcolor{cmtgray}{\#\ 3.\ 三条学习路径}}
\codeline{\textcolor{cmtgray}{\#\ ==============================}}
\codeblank{}
\codeline{\textcolor{cmtgray}{\#\ 路径一：理论完整路径（研究生课程\ /\ 系统学习）}}
\codeline{\textcolor{cmtgray}{\#\ \ \ 1\ \(\rightarrow\)\ 2\ \(\rightarrow\)\ 3\ \(\rightarrow\)\ 4\ \(\rightarrow\)\ 5\ \(\rightarrow\)\ 6\ \(\rightarrow\)\ 7\ \(\rightarrow\)\ 8\ \(\rightarrow\)\ 9}}
\codeline{\textcolor{cmtgray}{\#\ \ \ 重点：第2章描述逻辑与第5章推理机制逐节精读}}
\codeblank{}
\codeline{\textcolor{cmtgray}{\#\ 路径二：工程速成路径（在职工程师\ /\ 项目驱动）}}
\codeline{\textcolor{cmtgray}{\#\ \ \ 1\ \(\rightarrow\)\ 3（CQ与七步法）\(\rightarrow\)\ 4（OWL与SPARQL）\(\rightarrow\)\ 7（建库）\(\rightarrow\)\ 9（实战）}}
\codeline{\textcolor{cmtgray}{\#\ \ \ 跳过：Tableau算法细节、概率推理；用到时回查第5章}}
\codeblank{}
\codeline{\textcolor{cmtgray}{\#\ 路径三：AI应用路径（大模型应用开发者）}}
\codeline{\textcolor{cmtgray}{\#\ \ \ 1\ \(\rightarrow\)\ 8（先看本体能为LLM做什么）\(\rightarrow\)\ 4（补语言基础）}}
\codeline{\textcolor{cmtgray}{\#\ \ \ \ \ \(\rightarrow\)\ 7（GraphRAG的图谱底座）\(\rightarrow\)\ 5（推理校验）\(\rightarrow\)\ 9}}
\codeline{\textcolor{cmtgray}{\#\ \ \ 重点：第8章幻觉控制三层架构\ +\ 第7章SHACL校验}}
\codeblank{}
\codeline{\textcolor{cmtgray}{\#\ ==============================}}
\codeline{\textcolor{cmtgray}{\#\ 4.\ 各章配套代码速查}}
\codeline{\textcolor{cmtgray}{\#\ ==============================}}
\codeline{\textcolor{cmtgray}{\#\ 第2章\ \ examples/\ \ 概念与逻辑示例（纯文本推演）}}
\codeline{\textcolor{cmtgray}{\#\ 第3章\ \ examples/\ \ CQ设计、七步法走查、OntoClean标注}}
\codeline{\textcolor{cmtgray}{\#\ 第4章\ \ examples/\ \ .ttl/.rdf/.owl/.sparql\ 可直接在Protégé/Jena中使用}}
\codeline{\textcolor{cmtgray}{\#\ 第5章\ \ examples/\ \ SWRL规则、Tableau推演、时态与概率示例}}
\codeline{\textcolor{cmtgray}{\#\ 第6章\ \ examples/\ \ 四大领域应用的本体片段与规则}}
\codeline{\textcolor{cmtgray}{\#\ 第7章\ \ examples/\ \ 构建流水线、实体对齐、SHACL\ shapes（.ttl）}}
\codeline{\textcolor{cmtgray}{\#\ 第8章\ \ examples/\ \ 幻觉控制架构、GraphRAG、Text2SPARQL、Agent伪代码}}
\codeline{\textcolor{cmtgray}{\#\ 第9章\ \ src/\ \ \ \ \ \ \ 可运行的\ Java（Jena）与\ Python（owlready2）源码}}
\end{codebox}
